% Part: computability
% Chapter: recursive-functions
% Section: pr-relations

\documentclass[../../../include/open-logic-section]{subfiles}

\begin{document}

\olfileid{cmp}{rec}{prr}
\olsection{روابط بازگشتی اولیه}


\begin{defn}
رابطهٔ~$R(\vec x)$ را بازگشتی اولیه می‌نامیم اگر تابع مشخصهٔ آن،
\[
\Char{R}(\vec x) = \left\{
  \begin{array}{ll}
  1 & \mbox{اگر $R(\vec x)$} \\
  0 & \mbox{در غیر این صورت}
  \end{array}
\right.
\]
بازگشتی اولیه باشد.
\end{defn}

به بیان دیگر، وقتی از رابطهٔ بازگشتی اولیه~$R(\vec x)$ سخن می‌گوییم،
منظور رابطه‌ای به صورت $\Char{R}(\vec x) = 1$ است که در آن
$\Char{R}$ تابعی بازگشتی اولیه است که برای هر ورودی یا~$1$ و یا~$0$
بازمی‌گرداند. برای نمونه، رابطهٔ~$\fn{IsZero}(x)$ که هرگاه و تنها
هرگاه $x = 0$ برقرار است، متناظر با تابع~$\Char{\fn{IsZero}}$ است که
با بازگشت اولیه چنین تعریف می‌شود:
\begin{align*}
\Char{\fn{IsZero}}(0) & = 1,\\
\Char{\fn{IsZero}}(x+1) & = 0.
\end{align*}

باید روشن باشد که می‌توان روابط را با دیگر توابع بازگشتی اولیه ترکیب
کرد. پس موارد زیر نیز بازگشتی اولیه‌اند:
\begin{enumerate}
\item رابطهٔ برابری~$x = y$ که با
$\fn{IsZero}(\left|x-y\right|)$ تعریف می‌شود؛
\item رابطهٔ کوچک‌تر یا مساوی~$x \leq y$ که با
$\fn{IsZero}(x \tsub y)$ تعریف می‌شود.
\end{enumerate}


\begin{prop}
  مجموعهٔ روابط بازگشتی اولیه نسبت به عملیات بولی بسته است؛ یعنی اگر
  $P(\vec x)$ و~$Q(\vec x)$ بازگشتی اولیه باشند، روابط زیر نیز چنین‌اند:
  \begin{enumerate}
  \item $\lnot P(\vec x)$؛
  \item $P(\vec x) \land Q(\vec x)$؛
  \item $P(\vec x) \lor Q(\vec x)$؛
  \item $P(\vec x) \lif Q(\vec x)$.
  \end{enumerate}
\end{prop}

\begin{proof}
  فرض کنید $P(\vec x)$ و~$Q(\vec x)$ بازگشتی اولیه‌اند؛ یعنی توابع
  مشخصهٔ آن‌ها، $\Char{P}$ و~$\Char{Q}$، بازگشتی اولیه‌اند. باید نشان
  دهیم توابع مشخصهٔ $\lnot P(\vec x)$ و جز آن نیز بازگشتی اولیه‌اند.
  \[
  \Char{\lnot P}(\vec x) = \begin{cases}
    0 & \text{اگر $\Char{P}(\vec x) = 1$}\\
    1 & \text{در غیر این صورت}
  \end{cases}
  \]
  می‌توانیم $\Char{\lnot P}(\vec x)$ را به صورت
  $1 \tsub \Char{P}(\vec x)$ تعریف کنیم.
  \[
  \Char{P \land Q}(\vec x) = \begin{cases}
    1 & \text{اگر $\Char{P}(\vec x) = \Char{Q}(\vec x) = 1$}\\
    0 & \text{در غیر این صورت}
  \end{cases}
  \]
  می‌توانیم $\Char{P \land Q}(\vec x)$ را به صورت
  $\Char{P}(\vec x) \cdot \Char{Q}(\vec x)$ یا
  $\fn{min}(\Char{P}(\vec x), \Char{Q}(\vec x))$ تعریف کنیم. به همین
  ترتیب،
  \begin{align*}
    \Char{P \lor Q}(\vec x) & = \fn{max}(\Char{P}(\vec x), \Char{Q}(\vec x)) \text{ و}\\
    \Char{P \lif Q}(\vec x) & = \fn{max}(1 \tsub
  \Char{P}(\vec x), \Char{Q}(\vec x)).
  \end{align*}
\end{proof}

\begin{prop}
  مجموعهٔ روابط بازگشتی اولیه نسبت به سوربندی کراندار بسته است؛ یعنی
  اگر $R(\vec x, z)$ رابطه‌ای بازگشتی اولیه باشد، روابط
  \begin{align*}
    & \bforall{z < y}{R(\vec x, z)} \text{ و}\\
    & \bexists{z < y}{R(\vec x, z)}
  \end{align*}
  نیز چنین‌اند. $\bforall{z < y}{R(\vec x, z)}$ برای~$\vec x$ و~$y$
  هرگاه و تنها هرگاه برقرار است که $R(\vec x,z)$ برای هر $z<y$
  برقرار باشد؛ و برای~$\bexists{z < y}{R(\vec x,z)}$ نیز به‌طور
  مشابه.
\end{prop}

\begin{proof}
  بنا بر قرارداد، $\bforall{z < 0}{R(\vec x, z)}$ را صادق می‌گیریم
  (به این دلیل ساده که هیچ $z<0$ای وجود ندارد) و
  $\bexists{z < 0}{R(\vec x, z)}$ را کاذب می‌گیریم. سور کلی کراندار
  درست مانند حاصل‌ضرب متناهی یا کمینهٔ تکرارشده عمل می‌کند؛ یعنی اگر
  $P(\vec x, y) \defiff \bforall{z < y}{R(\vec x, z)}$، می‌توان
  $\Char{P}(\vec x,y)$ را چنین تعریف کرد:
  \begin{align*}
    \Char{P}(\vec x, 0) & = 1\\
    \Char{P}(\vec x, y+1) & =
    \fn{min}(\Char{P}(\vec x, y), \Char{R}(\vec x, y)).
  \end{align*}
  سور وجودی کراندار را نیز می‌توان به‌طور مشابه با~$\fn{max}$ تعریف
  کرد. راه دیگر آن است که با استفاده از هم‌ارزی
  $\bexists{z < y}{R(\vec x, z)}
  \liff \lnot \bforall{z < y}{\lnot R(\vec x, z)}$، آن را از سور
  کلی کراندار تعریف کنیم. توجه کنید که، برای نمونه، سور کرانداری به
  صورت $\bexists{x \leq y}{\dots x\dots}$ با
  $\bexists{x < y+1}{\dots x \dots}$ هم‌ارز است.
\end{proof}

\begin{prob}
  نشان دهید رابطهٔ سه‌موضعی $x \equiv y \mod n$ (هم‌نهشتی به پیمانهٔ
  $n$) بازگشتی اولیه است.
\end{prob}

تابع بازگشتی اولیهٔ سودمند دیگری تابع شرطی~$\fn{cond}(x,y,z)$ است که
چنین تعریف می‌شود:
\begin{align*}
  \fn{cond}(x,y,z) & = \begin{cases}
  y & \text{اگر $x = 0$} \\
  z & \text{در غیر این صورت}.
\end{cases}
\intertext{این تابع به‌طور بازگشتی چنین تعریف می‌شود:}
\fn{cond}(0,y,z) & = y,\\
\fn{cond}(x+1,y,z) & = z.
\end{align*}
می‌توان از این تابع برای توجیه تعریف موردی توابع بازگشتی اولیه بر پایهٔ
روابط بازگشتی اولیه استفاده کرد:

\begin{prop}
اگر $g_0(\vec x)$، \dots، $g_m(\vec x)$ توابع بازگشتی اولیه و
$R_0(\vec x)$، \dots، $R_{m-1}(\vec x)$ روابط بازگشتی اولیه باشند،
تابع~$f$ که با
\[
f(\vec x) = \begin{cases}
    g_0(\vec x) & \text{اگر $R_0(\vec{x})$} \\
    g_1(\vec x) & \text{اگر $R_1(\vec{x})$ و نه $R_0(\vec{x})$} \\
    \vdots & \\
    g_{m-1}(\vec x) & \text{اگر $R_{m-1}(\vec{x})$ و هیچ‌یک از
      حالت‌های پیشین برقرار نباشد}
    \\
    g_m(\vec x) & \mbox{در غیر این صورت}
\end{cases}
\]
تعریف می‌شود نیز بازگشتی اولیه است.
\end{prop}

\begin{proof}
  اگر $m = 1$، این همان تابعی است که با
  \[
  f(\vec x) = \fn{cond}(\Char{\lnot R_0}(\vec x),g_0(\vec x),g_1(\vec
  x))
  \]
  تعریف می‌شود. برای $m>1$ کافی است تعریف‌هایی از این صورت را با هم
  ترکیب کنیم.
\end{proof}
\end{document}
